% !TeX root = ../papers/nc-p3.tex
\section{対合が交換する四点と七次元成分の全次数判定}
\label{sec:involution-quartet-component}
四本の四周期点列が各位置で独立とし、その点座標を $0,1,2,3$ の座標点とする。
従って同じラベルが巡回して隣接する係数は零である。残る84係数が非零である開集合を取る。
対合が両対 $I=\{0,1\},J=\{2,3\}$ の各々を交換するとし、各基底の倍率を合わせて
四位置で $(01)(23)$ とする。テンソルを対合の正固有空間に置く。
四つの関係空間で対合の正・負部分の次元は $4,2$ とする。

\begin{lemma}\label{lem:four-track-opposite-minors}
任意の反対位置で、一方に $I$ の点、他方に $J$ の点を入れた二文字縮約が全て階数1とする。
階数1関係がなく、四点列での左右の関係評価が各々階数3なら、窓は標準Koszul左右AS正則である。
\end{lemma}
\begin{proof}
各位置で一対の点の消滅二平面を $U$ とする。その対の一方を $p$、他方を $q$ と書く。
四次元の正関係を $p$ で評価すると後続点の消滅超平面に入る。
従って正関係の一つ $r\ne0$ が $p$ を消し、対合で $q$ も消す。
点 $p^-$ で縮約した三テンソルの末尾縮約は $K(p)$ 全体を張るので
$r=w(p^-,-,-,l)$ と書ける。$w(p^-,q,-,-)$ の階数は2である。
実際、$K(p)$ の $q$ 評価は $r$ を消し、点評価階数3・二文字の完全対合によりこの縮約は非零の関係である。
次位置の他方の対のどちらの点 $b$ に対しても $w(p^-,-,b,-)$ は階数1である。
その $q$ 行は零でない。四点が座標点の表示で許された成分が一つ以上あり、全て非零だからである。
$l$ はこの行を消すので行列全体を消す。よって $r$ の右支持は他方の対の消滅二平面であり、
両対を交互に取ると定理\ref{thm:recurrent-point-line-flags}の条件2になる。

一歩の後続は $U'=\operatorname{supp}_{\rm left}w(p^-,q,-,-)$ で $q$ の後続点を消す。
二位置後の他方の対の点 $c$ では $w(q,-,c,-)$ が階数1で、その $p^-$ での非零列は $U'$ に入る。
従って $U'^\perp$ の任意の元 $s$ は $w(q,s,c,-)=0$ を満たす。
$s$ を $q$ の後続とは独立に取れば、二歩の後続関係は非零で階数2、左支持は二位置後の他方の対の消滅二平面に等しい。
これで条件1も得る。定理の証明で評価下界を使うのは指定点列上での三生成表示、前後の核の一意性、
および最後に単純を解消する一つの点列だけであり、全てこの四点列から取れる。
\end{proof}

\begin{theorem}\label{thm:involution-quartet-component}
反対階数16の点 $t$（命題\ref{prop:klein-rational-full-opposite}）を通る、
対合付きの七次元の局所成分がある。一般点は上の四点と階数1縮約の表示を持つ。
この成分の基底変更による像の閉包では、二隣接階数6、一文字全階数、階数1関係の不在、
全位置の左右評価下界3を満たす全ての閉点が標準Koszul左右AS正則である。
特に五次元楕円補完の既約族全体の閉包にこの含意が成立する。
\end{theorem}
\begin{proof}[四点表示の局所モデル]
同時に1を加える対合で残った42係数を使う。$w_{ijkl}$ は第一ラベルが0または2の方で表す。
左二文字・右二文字をいずれも $01,23,02,03,20,21$ の順の六軌道で取る。
開始位置 $s=0,1$、符号 $\epsilon=\pm1$ ごとに
$w_{ijkl}+\epsilon w_{ij,k+1,l+1}$ の六行列を $M_s^\epsilon$ と書く
（加算はビット排他的論理和、$s=1$ では四位置を巡回する）。
前半二行・列を一組、後ろを二行・列ずつ二組とする。
反対階数1の条件は前後に残る二ラベルの二行列式が零であることなので八つの二項式になる。
次の表の記号を使う。
\[
\begin{array}{c|rrrrrrrrrr}
 &P&Q&A&B&C&D&E&F&G&H\\ \hline
 w&0101&2323&0123&0132&2301&2310&0231&0321&2013&2103\\
\end{array}\qquad
\begin{array}{c|rrrrrrrr}
 &i&j&k&l&m&n&o&r\\ \hline
 w&0102&0103&0121&0131&0201&0301&2010&2101
\end{array}
\]
八つの二項式は別々の成分 $2023,2123,0323,0232,2320,2321,2313,2303$ を消す。
例えば $w_{0102}w_{2023}=w_{0132}w_{2013}$ である。
全体倍率と、各位置の二対の相対倍率で $P=Q=A=E=l=1$ とできる。
実際、全体を $1/P$ 倍し、位置 $2,3,1,0$ の順に相対倍率を決めればよい。

七変数 $x,X,u,U,B,F,k$ を取り、次の値を入れる。
\[
\begin{aligned}
 C&=(1-x)(1-X),&D&=xX/B,&G&=uU,&H&=(1-u)(1-U)/F,\\
 i&=Bu,&j&=(1-u)/k,&m&=x,&n&=F(1-x)/k,\\
 o&=XU/B,&r&=k(1-X)(1-U)/F .
\end{aligned}
\]
これらの変数は表示の係数から復元できる。実際に $x=lm/E$, $X=io/G$,
$u=il/B$, $U=mo/D$ である。
両 $M_s^-$ の先頭二行列を可逆とし、そのSchur補を取る。後ろ四行列の二つの交差ブロックは
この代入で恒等的に零となる。対角二ブロックの二切断合わせて16式は、残った交互配置
（例えば $0202,0213$）の16成分について定数係数の一次式で、係数階数は12である。
辞書式の自由座標を $w_{0312},w_{0313},w_{2130},w_{2131}$ とする。四つの整合条件を $f_j$ と書く。
第一式から単元を除けば $U$ について一次であり、第二式は $k^2$ について一次である。
これらを順に解くと残り二式は恒等的に零となる。従って
$x,X,u,B,F$ は独立であり、$k$ の二次拡大を許せば負切断の全式を満たす。

次に $M_s^+$ の先頭二行列のSchur補は二つの対角二行列からなり、交差ブロックは零である。
一方の対角行列には最初の自由二座標 $v,w$ だけ、他方には残る二座標だけが現れる。
各々の対角は $2v$ に定数を、非対角は $2w$ に定数を加えたものである。
第二切断では両組が交換され、対応する二行列式は二次項、一次項、定数項まで一致する。
従って両符号の必要な階数を得るため、各対の二座標についてちょうど一つずつのアフィン二次曲線を取ればよい
（各二行列が零になる閉集合を避ける）。これが七次元の表示である。
\path{computations/verify_involution_quartet_component.py}は上の二つの一次消去、四ブロックの
Schur補、二次曲線の一致を五変数有理関数体とその二次拡大で検算する。
この消去はこの表示を構成するためのもので、他の根の組合せを全てこの表示に含める必要はない。
\end{proof}

\begin{proof}[元の階数軌跡への支配と特殊化]
点 $t$ の素数2657での二十点のうち、位置0の
\[
 (1,2197,1186,92),\quad(1,2197,1471,2565),\quad
 (1,1667,2322,492),\quad(1,1667,335,2165)
\]
を順に取る。各対は高位二座標の符号を変える対合で交換される。
前節までの計算で点を四位置に移すと元に戻る。
これらを座標枠にすると84係数は全て非零で、八二項式を満たす。
再現計算では全42係数を上記表示から復元し、$U,k^2$ の二主係数と全ての使用分母が単元であることも確認する。

この剰余点の近傍で、まだ八二項式を課さずに二つの隣接切断の階数6条件を微分する。
各切断の左核と右核の間への微分の射影を、42座標について縦に並べた行列の階数は30である。
従って $\mathbb Z_{(2657)}$ 上でこの30個の独立な式だけを課した周囲は剰余点で相対次元12の滑らかな空間である。
上で構成した部分はそこに、五つの倍率を戻した相対次元12で入る。
次元の下界はこの点自身で成立する。五独立座標と二次根は主係数可逆の開部分にあり、
残る四変数には二式しかない。係数から全変数を復元する局所逆もある。
特に付加条件を課した局所モデルの絶対次元は少なくとも13で、上の正則局所環と等しい。
正則局所環は整域なのでこの中に同次元の真の閉部分はなく、階数軌跡と表示は局所的に一致する。
この議論は有限体の標本から標数零の恒等式を推測するものではない。
四点を忘れたとき一般の選択は有限で、枠を固定した後の連続な同一視は上の五倍率だけである。
点配置の安定群も標本で有限なので一般点で有限であり、基底変更を除いた次元は7となる。

五次元族の標本 $t$ は標数零から同じ剰余点に特殊化する。
単純な点配置を局所エタールに番号付けし、上の二点と対合による二つの像を取れる。
四移動は単純な点の番号の置換で、剰余点で恒等であり、点列の四周期性が保たれる。
標数零の一般点の点対応と四周期性は定理\ref{thm:exact-two-configurations}の証明からも従う。
従って五次元族の一般点のこの表示が上の開近傍に入り、既約性でその全閉包が七次元成分の
基底変更像の閉包に含まれる。一般点の評価下界・階数1関係の不在はそれらが成立する $t$ から開性で従う。
補題\ref{lem:four-track-opposite-minors}の巡回旗と定理\ref{thm:recurrent-point-line-flags}の固有性で、
表示が消える境界でも定理に述べた全閉点の正則性を得る。
\end{proof}

\begin{corollary}\label{cor:full-thin-finite-graph-regularity}
全成分非零の薄い二色を持つ平行切断の五次元主閉包では、純関係がなく、
四辺が有限であれば、一文字全階数・隣接階数6の候補は全て標準Koszul左右AS正則である。
特に混合交換表示でも、以前の四中心式で認証した開集合の外までこの結論が通用する。
\end{corollary}
\begin{proof}
片側の評価が階数2以下の点があれば、その点の相手が少なくとも直線をなし有限性に反する。
従って左右評価の下界を満たす。
\end{proof}
